\section{Reconstruction of the published span}
\label{app:b10}

\paragraph{The problem.}
The published degree-ten structures are given as explicit expressions. To compare
their span with the product subspace over $\Q$ --- rather than modulo a prime ---
their coordinate vectors must be lifted to $\Q$.

\paragraph{CRT and rational reconstruction.}
Each published structure was evaluated at $\bpFitPrimes$ fitting primes. The
residues were combined by the Chinese remainder theorem into a residue modulo the
product, and a rational with bounded numerator and denominator was recovered by
rational reconstruction~\cite{Wang:1982}. All $\dimBten$ published structures
lifted; none failed.

\paragraph{Held-out validation.}
The reconstructed rationals were reduced modulo the held-out prime
$\bpHoldoutPrime$, which took no part in the fit, and compared against a direct
evaluation there. There were no mismatches. This is the step that distinguishes a
reconstruction from a guess: rational reconstruction always returns
\emph{something}, and only a prime outside the fit can tell you whether it
returned the right thing.

\paragraph{Exact intersection.}
With the lifts in hand,
\begin{equation}
\dim_\Q \Bten = \dimBten,\quad
\dim_\Q \Pten = \dimPten,\quad
\dim_\Q(\Bten + \Pten) = \dimBplusP,
\end{equation}
whence $\dim_\Q(\Bten \cap \Pten) = \dimBcapP$ by inclusion--exclusion.

\paragraph{The explicit generator.}
The one-dimensional intersection is generated by
\begin{equation}
-55\,I^{(1)} - 252\,I^{(2)} + 4200\,I^{(5)} + 5040\,I^{(6)}
 \;=\; -1040\,I^{(4)}_1 I^{(6)}_1,
\end{equation}
an identity between integer combinations, verified at the fresh prime
$\bpFreshPrime$ on $\bpFreshSamples$ independently drawn samples that took no
part in either the fit or the holdout.

\paragraph{Normalisation, and a check that failed.}
The verification initially failed at the fresh prime on all six samples. The
cause was in the checking code: the two sides of the identity were each being
reduced by their own greatest common divisor before comparison, which destroys
precisely the relative scale the identity asserts. Recomputed on a single common
scale, the identity holds on all six samples.

We record the failure rather than only the eventual success. A check that fails,
is adjusted, and then passes is worthless unless the adjustment is stated, because
the adjustment is where the error would hide. The adjustment here was to stop
normalising the two sides independently; no coefficient was altered to make the
identity fit.

\paragraph{Why the two decompositions differ.}
The published structures were presented as explicit expressions, not as a basis
with specified relations, and there is no reason a set chosen for explicitness
should avoid the product sector. That it meets it in exactly one dimension is a
structural observation about the published choice, and the direction of the
finding --- that $\Bten$ has one dimension \emph{fewer} of primitive content than
a naive reading would suggest --- can only weaken a claim about the published
span, never strengthen one. No result in this paper depends on the equality.
